% !PRINTED 308
\weberpage{308}
%<<CONT
\[
\left|\begin{array}{c|cc|cc|cc|cc}
x-0 & PV, & PV & PP, & PV & VP, & VP & VV, & VP \\ \hline
x+0 & PP, & PP & PV, & PP & VV, & VV & VP, & VV
\end{array}\right|
\text{.}
\]
In the first case two $PP$ are \gesperrt{gained}, in the second and fourth cases no change occurs, in the third case two $VP$ are \gesperrt{lost}.

But with this the theorems II, III are proved.

The assumptions 2., 3., that no two adjacent ones of the functions $f_\nu$ and $F_\nu$ shall vanish for the same $x$, we may now give up. For if we hold fast to the first assumption, that no multiple roots are present, we can impart to the coefficients $a_\nu$ changes so small that, first, the assumptions 2., 3. are satisfied, and secondly, the roots lying between $\alpha$ and $\beta$ change so little that none of them leaves the interval. Since the roots occur only singly, real roots cannot in this process pass into imaginary ones either.

Whether the theorem, with correct counting of the multiple roots, remains valid also in the case of multiple roots may be left undecided.

\begin{center}
\S.\,101.

The Cartesian Theorem.
\end{center}

The \gesperrt{Cartesian} theorem, which is also named after \gesperrt{Harriot}\footnote{Harriot, Artis analyticae praxis ad aequationes algebraicas resolvendas 1631. Gauss' Werke, Bd. III, S. 67. Lagrange, l. c.}, gives a rule for estimating the number of positive roots. We can regard it simply as a special case of \gesperrt{Budan}'s theorem, by taking $\alpha = 0$ and $\beta = \infty$, or at least so large that the functions
\weq{1}{f(x), f'(x), f''(x) \ldots f^{(n)}(x)}
all have the same sign for $x > \beta$. This sign agrees with the sign of the coefficient of the highest power of $x$ in $f(x)$, and may be assumed positive.

If
\[
f(x) = a_0 x^n + a_1 x^{n-1} + a_2 x^{n-2} + \cdots + a_{n-1} x + a_n,
\]
%CONT>>
